% ============================================================================
% IEEE Conference Paper: Adaptive-MMRAG
% Author: Abhinay Sivva
% Format: IEEEtran conference style
% ============================================================================

\documentclass[conference]{IEEEtran}

% ---- Packages ---------------------------------------------------------------
\usepackage{cite}
\usepackage{amsmath,amssymb,amsfonts}
\usepackage{algorithmic}
\usepackage{algorithm}
\usepackage{graphicx}
\usepackage{textcomp}
\usepackage{xcolor}
\usepackage{url}
\usepackage{booktabs}
\usepackage{multirow}
\usepackage{array}
\usepackage{balance}

\def\BibTeX{{\rm B\kern-.05em{\sc i\kern-.025em b}\kern-.08em
    T\kern-.1667em\lower.7ex\hbox{E}\kern-.125emX}}

% ============================================================================
\begin{document}

\title{Adaptive-MMRAG: A Multi-Agent, Resource-Aware, Offline\\
Multimodal Retrieval-Augmented Generation Framework\\
with Dual-Stage Evidence Verification}

\author{
  \IEEEauthorblockN{Abhinay Sivva\IEEEauthorrefmark{1},
                    Vankadaru Hari Kiran\IEEEauthorrefmark{1},
                    Adithya Gaddam\IEEEauthorrefmark{1},
                    Yashwanth\IEEEauthorrefmark{1}}
  \IEEEauthorblockA{\IEEEauthorrefmark{1}%
    Department of Computer Science and Engineering,\\
    Amrita Vishwa Vidhyapeetam, Amaravati Campus,\\
    Andhra Pradesh, India\\
    \{abhinay, harikiran, adithya, yashwanth\}@gmail.com
  }
}

\maketitle

% ============================================================================
\begin{abstract}
Existing offline Retrieval-Augmented Generation (RAG) systems adopt static,
fixed-policy retrieval pipelines that map every query through an identical
retrieval depth, single hardcoded language model, and produce answers without
verifying whether generated claims are actually grounded in the retrieved
evidence. On resource-constrained edge devices, this design causes high
hallucination rates (up to 28\%), excessive memory pressure (8.4\,GB peak),
and sub-optimal query latency. We present \textit{Adaptive-MMRAG}, a fully
offline, privacy-preserving multimodal RAG framework that introduces three
principal contributions: (i) a six-agent adaptive decision pipeline that
dynamically classifies query intent, complexity, and target modalities, then
allocates retrieval budgets accordingly; (ii) a deterministic,
resource-aware multi-model router that monitors real-time RAM and CPU
utilization via \texttt{psutil} and selects between lightweight and
high-capacity local LLMs to prevent out-of-memory failures while preserving
reasoning depth; and (iii) a dual-stage evidence verification mechanism
coupling pre-generation retrieval confidence scoring with post-generation
factual claim auditing and formal abstention enforcement. On an internal
ablation benchmark, Adaptive-MMRAG reduces hallucination from 28.0\% to
2.5\%, raises faithfulness from 0.68 to 0.96, achieves 94.2\% abstention
accuracy on unanswerable queries, and cuts peak RAM consumption from 8.4\,GB
to 4.7\,GB relative to the static baseline. The complete system—ingesting
PDF, DOCX, scanned documents (via Tesseract OCR), images, and audio—operates
with zero external API calls.
\end{abstract}

\begin{IEEEkeywords}
Retrieval-Augmented Generation, multimodal RAG, offline inference, edge AI,
hallucination mitigation, evidence verification, multi-agent systems,
resource-aware routing, adaptive retrieval, privacy-preserving NLP
\end{IEEEkeywords}

% ============================================================================
\section{Introduction}\label{sec:intro}

Retrieval-Augmented Generation (RAG) \cite{lewis2020rag} has become the
dominant strategy for grounding large language model (LLM) responses in
factual, up-to-date document corpora. The core idea—retrieve relevant
passages and condition generation on them—substantially reduces hallucination
compared to parametric-only generation. However, several structural
limitations constrain production deployments, particularly in
privacy-sensitive or bandwidth-restricted settings.

\textbf{Static retrieval depth.} Conventional RAG systems retrieve a fixed
number of top-$K$ documents regardless of whether the query requires a
single sentence lookup or a multi-document synthesis. This blunt approach
wastes compute on simple queries and under-retrieves on complex ones.

\textbf{Single-model inference.} Most offline RAG systems hardcode a single
local LLM. On edge hardware with limited RAM (4–8\,GB), loading a
high-capacity model to answer a trivial factual question is wasteful; yet
routing every complex reasoning request to a lightweight model degrades
answer quality significantly.

\textbf{No post-generation verification.} Pre-retrieval confidence scores
indicate whether \textit{retrieved passages} are relevant, but do not prevent
the generation model from fabricating claims unsupported by those passages.
Without a post-generation audit, systems confidently emit hallucinated
answers.

\textbf{Privacy constraints.} Enterprise, medical, and legal deployments
often prohibit sending documents to cloud APIs. Existing high-quality
multimodal RAG solutions—including GPT-4V \cite{openai2023gpt4}, Gemini
\cite{google2023gemini}, and LLaVA-based cloud services—require internet
connectivity and transmit data externally, violating local data sovereignty
requirements.

We address these limitations with \textit{Adaptive-MMRAG}, an entirely
offline framework built around six specialized agents, a deterministic
resource-aware model router, and a dual-stage evidence verifier. The system
supports heterogeneous document types—text-based PDFs, DOCX files with
embedded images, scanned documents (via Tesseract OCR \cite{smith2007tesseract}),
standalone images, and audio transcripts—all processed locally using
Whisper \cite{radford2023whisper}, CLIP \cite{radford2021clip},
FAISS \cite{johnson2019billion}, and Ollama-hosted language models.

The paper makes the following empirically validated contributions:
\begin{enumerate}
  \item A six-agent adaptive pipeline that classifies query intent and
        complexity in $<\!2$\,ms without invoking an LLM, then selects
        retrieval strategy, modality scope, and top-$K$ budget accordingly.
  \item A dual-stage evidence verification mechanism: (a) pre-generation
        retrieval confidence scoring triggers adaptive re-retrieval when
        confidence $<\!0.50$, and (b) post-generation claim auditing enforces
        formal abstention when faithfulness $<\!0.25$ or attaches grounding
        warnings when faithfulness $\in [0.25, 0.60)$.
  \item A deterministic resource-aware model router that reads real-time
        system RAM via \texttt{psutil} and routes to a Tier~1 lightweight
        model (\texttt{qwen2.5:3b}) when available RAM $<\!6.0$\,GB, or to a
        Tier~2 reasoning model (\texttt{llama3.1:8b}) for high-complexity
        queries with sufficient memory, reducing hallucination and peak RAM
        simultaneously.
  \item A version-aware metadata schema with SHA-256 content hashing that
        prevents semantic similarity collisions between concurrent document
        revisions.
  \item A bidirectional offline audio subsystem that couples local automatic
        speech recognition (Whisper) for voice queries with an on-device
        text-to-speech engine (\texttt{pyttsx3}) synthesizing natural speech
        answers locally in $<\!0.5$\,s with zero extra GPU or RAM footprint.
\end{enumerate}

% ============================================================================
\section{Related Work}\label{sec:related}

\subsection{Retrieval-Augmented Generation}
Lewis et al.\ \cite{lewis2020rag} formalized the RAG paradigm, demonstrating
that conditioning generation on retrieved passages substantially outperforms
closed-book generation on knowledge-intensive tasks. Subsequent work
including FiD \cite{izacard2021leveraging} and Atlas \cite{izacard2022atlas}
extended dense retrieval to multi-document fusion. These systems operate on
text corpora and assume cloud-scale compute, neither addressing multimodal
inputs nor offline constraints.

\subsection{Multimodal Retrieval}
CLIP \cite{radford2021clip} established a shared vision-language embedding
space enabling cross-modal retrieval. FAISS \cite{johnson2019billion} provides
efficient billion-scale nearest-neighbor search. BM25 \cite{robertson2009probabilistic}
remains a competitive lexical baseline; Reciprocal Rank Fusion (RRF)
\cite{cormack2009reciprocal} fuses dense and lexical ranked lists without
score normalization overhead. Our system combines all three in a hybrid
retrieval pipeline with Cross-Encoder reranking \cite{nogueira2019passage}.

\subsection{Hallucination in RAG}
Hallucination in RAG—where the model generates claims absent from or
contradicted by retrieved passages—is actively studied \cite{huang2023survey,
ji2023survey}. RAGAS \cite{es2023ragas} introduced automated faithfulness
and answer relevancy metrics. Self-RAG \cite{asai2024selfrag} trains models
to emit reflection tokens to decide whether retrieval is needed and whether
generated statements are supported; this approach requires fine-tuning,
which is infeasible in a fully offline, no-GPU-fine-tuning setting.
ActiveRAG \cite{wang2024activerag} explores query-driven adaptive retrieval
but does not provide resource-aware model selection or post-generation
abstention. Our approach enforces hallucination mitigation without
fine-tuning, relying instead on a verifier agent applied after generation.

\subsection{Offline and Edge AI}
Offline LLM deployment using frameworks such as Ollama \cite{ollama2023},
llama.cpp \cite{llamacpp}, and GPT4All \cite{gpt4all} enables local inference
without cloud connectivity. However, these frameworks provide only inference
backends; they do not address adaptive retrieval, multi-model routing, or
post-generation verification. Our work builds on Ollama as the inference
backend while contributing the full adaptive pipeline around it.

\subsection{Adaptive Retrieval}
Adaptive retrieval selects retrieval depth or strategy based on query
characteristics \cite{mallen2023not, asai2024selfrag}. Existing adaptive
methods are predominantly designed for single-modality text corpora and
require either fine-tuned models or cloud API access for complexity
classification. Our query analysis agent performs deterministic, heuristic
classification in under 2\,ms using regular expression pattern matching and
keyword features—no LLM inference required for the routing decision.

% ============================================================================
\section{System Architecture and Methodology}\label{sec:methodology}

\subsection{Overview}

Adaptive-MMRAG is a six-agent pipeline operating entirely on local hardware.
Figure~\ref{fig:pipeline} illustrates the end-to-end flow. At query time, the
pipeline executes the following stages:

\begin{enumerate}
  \item Query Analysis (Agent~1): intent, complexity, modality classification.
  \item Retrieval Planning (Agent~2): strategy, top-$K$, reranker selection.
  \item Multimodal Hybrid Search (Agent~3): FAISS + BM25 + RRF + Cross-Encoder.
  \item Pre-Generation Evidence Verification (Agent~4): confidence scoring
        and adaptive re-retrieval trigger.
  \item Resource-Aware Model Routing: deterministic tier selection.
  \item Grounded Response Generation (Agent~5) + Post-Generation Claim
        Verification and Abstention (Agent~6).
\end{enumerate}

\begin{figure}[t]
\centering
\fbox{\parbox{0.95\columnwidth}{\centering\small
\textbf{User Query}\\[3pt]
$\downarrow$\\[2pt]
\textbf{Agent 1: Query Analysis} (intent, complexity, modality)\\[2pt]
$\downarrow$\\[2pt]
\textbf{Agent 2: Retrieval Planner} (strategy, top-$K$, rerank flag)\\[2pt]
$\downarrow$\\[2pt]
\textbf{Agent 3: Multimodal Hybrid Search}\\
\small{FAISS (512-d CLIP) + BM25Okapi + RRF + Cross-Encoder}\\[2pt]
$\downarrow$\\[2pt]
\textbf{Agent 4: Evidence Verifier} ($S \in [0,1]$)\\[2pt]
$S < 0.50$ $\Rightarrow$ \textit{Adaptive Re-Retrieval Loop}
\ $\uparrow$ (max 2 iterations)\\[2pt]
$S \geq 0.50$ $\downarrow$\\[2pt]
\textbf{Resource-Aware ModelRouter} (psutil RAM check)\\
RAM $<$ 6\,GB $\Rightarrow$ Tier~1 \quad RAM $\geq$ 6\,GB, complex $\Rightarrow$ Tier~2\\[2pt]
$\downarrow$\\[2pt]
\textbf{Agent 5: Grounded Response Generation} (streaming)\\[2pt]
$\downarrow$\\[2pt]
\textbf{Agent 6: Post-Generation Claim Verifier}\\
Faithfulness $\geq 0.60$: answer + citations\\
$[0.25, 0.60)$: grounding warning prefix\\
$< 0.25$: formal abstention enforced
}}
\caption{Adaptive-MMRAG end-to-end pipeline. Arrows with feedback loops
indicate the adaptive re-retrieval cycle controlled by Agent~4.}
\label{fig:pipeline}
\end{figure}

\subsection{Document Ingestion and Multimodal Indexing}
\label{subsec:indexing}

The offline ingestion pipeline processes five document modalities through
parallel workers:

\textbf{Text PDFs and DOCX.} Text extraction uses the \texttt{unstructured}
library \cite{unstructuredio} with \texttt{pdfplumber} fallback. Embedded
table cells are preserved as structured text chunks.

\textbf{Scanned and Flat PDFs.} When no selectable text is detected, the
system falls back to a \texttt{pypdfium2} rendering pipeline followed by
Tesseract OCR \cite{smith2007tesseract} (version 5.x), which processes each
page as a rendered image and extracts text using trained neural network
models (LSTM-based).

\textbf{Embedded Images.} Images extracted from PDFs and DOCX documents are
processed through \texttt{pytesseract} for OCR text extraction and
simultaneously encoded by the CLIP \cite{radford2021clip} vision encoder
(\texttt{clip-ViT-B-32}) into 512-dimensional vectors.

\textbf{Audio.} Audio files (.mp3, .wav) are processed by OpenAI
Whisper \cite{radford2023whisper} running locally (base model, int8 quantized)
to produce transcripts. Transcripts are chunked and embedded identically to
text.

All chunks are embedded using a unified \texttt{SentenceTransformer}
\texttt{clip-ViT-B-32} model producing 512-dimensional L2-normalized vectors
stored in a FAISS flat index (\texttt{IndexFlatIP}) \cite{johnson2019billion}.
To eliminate repeated disk I/O latency, the embedding generator and cross-encoder
reranker instances are preserved in memory as resident singletons across requests.
Embedding inference is vectorized across chunk batches ($B=32$), reducing
vectorization latency by over $80\%$.

\textbf{Digital-First High-Speed Extraction.} Ingestion adopts a digital-first
parsing pipeline: digital text and tables are extracted natively in $<100$\,ms via
\texttt{pdfplumber}. When character count falls below a minimum threshold ($<50$ characters),
indicating a scanned artifact, the system automatically transitions to the
\texttt{pypdfium2} and Tesseract OCR engine. For single- or dual-document uploads,
processing executes directly in-process to avoid the multi-second process-pool cold-start
penalty on Windows systems.

\textbf{Automated NLP Profiling and Keyphrase Extraction.} Concurrently with vector indexing,
an offline Natural Language Processing (NLP) summarization engine generates document-level
metadata profiles. Salient domain keyphrases are extracted via position-weighted term-frequency
analysis and multi-word collocation scoring, while executive summaries are generated through
sentence centrality scoring and salience density ranking in $<15$\,ms without requiring LLM
generation. These keyphrases serve as interactive query suggestions on the user interface,
enabling zero-latency exploration of indexed documents.

\textbf{Version-Aware Metadata Schema.} Every chunk carries extended
metadata: \texttt{chunk\_id}, \texttt{source}, \texttt{page\_number},
\texttt{type}, \texttt{document\_id}, \texttt{version} (parsed from filename
pattern \texttt{\_v$N$}), \texttt{content\_hash} (SHA-256), and
\texttt{timestamp}. During retrieval, a \texttt{prioritize\_version\_aware\_candidates()}
function sorts candidates of the same logical document by descending version
number, preventing stale document revisions from displacing current ones.

\subsection{Agent 1: Query Analysis Agent}
\label{subsec:agent1}

The Query Analysis Agent classifies three dimensions of each incoming query
using deterministic heuristics (no LLM inference):

\textbf{Intent} (9 classes): \textit{factual}, \textit{summarization},
\textit{comparison}, \textit{visual\_question}, \textit{audio\_question},
\textit{cross\_modal}, \textit{reasoning}, \textit{document\_overview},
\textit{keyword\_lookup}. Classification uses keyword pattern matching on
interrogative tokens, presence of visual/audio triggers (``image'',
``diagram'', ``transcript''), and comparison markers (``compare'',
``difference'', ``versus'').

\textbf{Complexity} (3 levels: LOW, MEDIUM, HIGH): LOW if the query matches
direct lookup patterns with $<\!5$ content tokens; HIGH if multi-hop
reasoning markers (``step by step'', ``explain why'', ``compare all'') or
overview requests (\texttt{OVERVIEW}) are present; MEDIUM otherwise.

\textbf{Target Modalities} (TEXT, IMAGE, AUDIO, TABLE, MULTIMODAL):
Modality scope is narrowed or expanded based on intent. A
\textit{visual\_question} intent restricts search to IMAGE and TABLE chunks;
\textit{audio\_question} restricts to AUDIO; \textit{cross\_modal} enables
full multimodal fusion.

Total execution time for Agent~1: $<\!2$\,ms on the test hardware (16-core
CPU, no GPU).

\subsection{Agent 2: Retrieval Planner Agent}
\label{subsec:agent2}

The Retrieval Planner translates the query analysis into a concrete
$\texttt{RetrievalPlan}$ specifying:
\begin{itemize}
  \item \textbf{Top-$K$ budget}: $K=4$ (LOW), $K=5$ (MEDIUM), $K=8$ (HIGH),
        $K=12$ (OVERVIEW). During adaptive re-retrieval, $K$ is expanded by a
        factor of 1.4.
  \item \textbf{Reranker flag}: Cross-Encoder reranking enabled when
        complexity $\geq$ MEDIUM or re-retrieval is triggered.
  \item \textbf{Retrieval strategy tag}: one of
        \texttt{hybrid\_rerank}, \texttt{adaptive\_broadened\_hybrid\_rerank},
        \texttt{semantic\_only}, or \texttt{keyword\_focused}.
\end{itemize}

\subsection{Agent 3: Multimodal Hybrid Search}
\label{subsec:agent3}

Retrieval fuses three complementary signals:

\textbf{Dense Semantic Search.} The query is embedded into 512-d CLIP vector
space and compared via inner product (FAISS \texttt{IndexFlatIP}) to all
indexed chunk embeddings, yielding ranked candidate list $L_{\text{dense}}$.

\textbf{BM25 Lexical Search.} A \texttt{BM25Okapi} \cite{robertson2009probabilistic}
index over chunk text fields returns ranked list $L_{\text{bm25}}$.

\textbf{Reciprocal Rank Fusion (RRF).} For candidate $d$ appearing in
both lists, the fused score is:
\begin{equation}
  \text{RRF}(d) = \sum_{L \in \{L_{\text{dense}}, L_{\text{bm25}}\}}
  \frac{1}{k + r_L(d)}
  \label{eq:rrf}
\end{equation}
where $r_L(d)$ is the rank of $d$ in list $L$ and $k=60$ (standard
constant \cite{cormack2009reciprocal}).

\textbf{Cross-Encoder Reranking.} When enabled, the top-$N$ RRF candidates
are reranked by \texttt{cross-encoder/ms-marco-MiniLM-L-6-v2}
\cite{nogueira2019passage} producing a relevance logit per
(query, chunk) pair.

\subsection{Agent 4: Pre-Generation Evidence Verifier and Adaptive
Re-Retrieval}
\label{subsec:agent4}

After retrieval, Agent~4 computes a normalized confidence score
$S \in [0,1]$:
\begin{equation}
  S = \alpha \cdot \bar{r}_3 + \beta \cdot c_{\text{cover}} + \gamma \cdot c_{\text{kw}}
  \label{eq:confidence}
\end{equation}
where $\bar{r}_3$ is the mean relevance score of the top-3 retrieved chunks
(normalized to $[0,1]$), $c_{\text{cover}}$ is the fraction of distinct
modalities represented in the retrieved set relative to the requested
modalities, $c_{\text{kw}}$ is the query keyword coverage ratio over the
retrieved text, and $\alpha + \beta + \gamma = 1$ (empirically set to
$0.5, 0.3, 0.2$).

Evidence is classified as HIGH ($S \geq 0.75$), MEDIUM ($S \geq 0.50$), or
LOW ($S < 0.50$). When LOW and fewer than $\text{MAX\_ITERATIONS}=2$
re-retrieval cycles have elapsed, the system triggers an
\textbf{Adaptive Re-Retrieval}: top-$K$ is expanded by $\times\!1.4$,
Cross-Encoder reranking is forced on, modality scope broadens to MULTIMODAL,
and retrieval is re-executed. The agent logs the initial and final
$S$ scores to the per-query telemetry trail.

\subsection{Resource-Aware Model Router}
\label{subsec:router}

The \texttt{ModelRouter} reads real-time system telemetry:
\begin{equation}
  \text{RAM}_{\text{avail}} = \texttt{psutil.virtual\_memory().available} / 2^{30}
  \label{eq:ram}
\end{equation}
and applies the following deterministic routing policy (summarized in
Algorithm~\ref{alg:router}):

\begin{algorithm}[h]
\caption{Deterministic Resource-Aware Model Routing}
\label{alg:router}
\begin{algorithmic}[1]
\REQUIRE $\text{RAM}_{\text{avail}}$, query complexity $C$,
         evidence confidence $E$, threshold $\theta = 6.0$\,GB
\IF{$\text{RAM}_{\text{avail}} < \theta$}
  \STATE Select Tier~1 (\texttt{qwen2.5:3b}), reason: OOM prevention
\ELSIF{$C = \text{LOW}$ \AND $E \in \{\text{HIGH, MEDIUM}\}$}
  \STATE Select Tier~1, reason: low complexity, fast path
\ELSIF{$C \in \{\text{HIGH, OVERVIEW}\}$ \AND $\text{RAM}_{\text{avail}} \geq \theta$}
  \STATE Select Tier~2 (\texttt{llama3.1:8b}), reason: deep reasoning
\ELSE
  \STATE Select Tier~1, reason: standard query, conservative path
\ENDIF
\RETURN (selected\_model, tier, reason, $\text{RAM}_{\text{avail}}$)
\end{algorithmic}
\end{algorithm}

The routing decision and all hardware telemetry are logged per query. In
the evaluation dataset, 64\% of queries were routed to Tier~1 and 36\% to
Tier~2; 8\% were RAM-constrained fallbacks that overrode a Tier~2 decision.

\subsection{Agent 5: Grounded Response Generation}
\label{subsec:agent5}

The Response Generation Agent formulates a structured prompt injecting the
retrieved evidence chunks with source citations, then streams tokens via the
selected local Ollama model. Prompt engineering enforces strict verbatim
grounding constraints: the model is instructed to answer solely from the
provided context and explicitly cite sources by chunk identifier.

\subsection{Agent 6: Post-Generation Claim Verifier and Formal Abstention}
\label{subsec:agent6}

After generation, Agent~6 extracts discrete factual claims from the generated
answer using sentence boundary detection and regex-based claim filtering
(discarding meta-statements such as ``here is a summary of...''). For each
extracted claim $c_i$, support is assessed against the evidence corpus using
three complementary signals:

\begin{enumerate}
  \item \textbf{Keyword overlap}: Let $W_i$ be the set of content words
        (excluding stopwords) in $c_i$ and $W_E$ the vocabulary of the
        evidence corpus. The overlap ratio $\rho_i = |W_i \cap W_E| /
        |W_i|$.
  \item \textbf{Numeric/date fidelity}: All numeric tokens in $c_i$ must
        appear verbatim in the evidence corpus.
  \item \textbf{Cross-Encoder verification}: When a reranker instance is
        available, the (claim, evidence summary) pair is scored; claims with
        logit $<\!-2.0$ are flagged unsupported.
\end{enumerate}

Claim $c_i$ is considered supported if $\rho_i \geq 0.35$ AND numeric
fidelity holds AND Cross-Encoder check passes. The faithfulness score is:
\begin{equation}
  F = \frac{|\{i : c_i \text{ supported}\}|}{|\{c_i\}|}
  \label{eq:faithfulness}
\end{equation}

\textbf{Abstention policy:}
\begin{itemize}
  \item $F \geq 0.60$: answer passed verbatim with citations.
  \item $F \in [0.25, 0.60)$: grounding warning banner prepended.
  \item $F < 0.25$: full formal abstention enforced; answer replaced with
        the system message \textit{``No supporting information was found in
        the available documents.''}.
\end{itemize}

\subsection{Bidirectional Offline Audio Subsystem}
\label{subsec:audio_subsystem}

To provide full hands-free accessibility without breaching offline privacy
guarantees, Adaptive-MMRAG incorporates a dual-phase audio pipeline:

\textbf{Speech-to-Text Voice Querying.} Spoken user queries recorded via the
client interface are streamed to the backend and transcribed using a local,
quantized instance of OpenAI Whisper \cite{radford2023whisper} running on CPU.
The transcribed text is then fed directly into the Query Analysis Agent
(Agent~1) as standard text input.

\textbf{Offline Text-to-Speech (TTS) Synthesis.} To close the audio loop,
generated responses are synthesized into speech using an on-device TTS handler
(\texttt{pyttsx3}) interfacing with native OS speech engines (e.g., Microsoft
David and Zira voices). Before synthesis, a dedicated natural language
sanitizer strips markdown headers, citation tokens (\texttt{[Source:\dots]}),
code fragments, and HTML formatting blocks. This ensures conversational,
fluent audio generation in under 0.5\,seconds without requiring cloud TTS
APIs or allocating additional VRAM.

% ============================================================================
\section{Implementation Details}\label{sec:implementation}

\subsection{Technology Stack}
All components run locally on commodity hardware without GPU requirement:
FAISS \cite{johnson2019billion} (CPU flat index), Sentence-Transformers
\cite{reimers2019sentence} (\texttt{clip-ViT-B-32}), rank-bm25
\cite{rankbm25}, Tesseract OCR 5.x \cite{smith2007tesseract}, OpenAI Whisper
(base, local) \cite{radford2023whisper}, Ollama \cite{ollama2023} for local
LLM inference, Flask for the web API, and standard NumPy/SciPy/Pillow for
numeric and image operations. The full dependency list is distributed via
\texttt{requirements.txt}.

\subsection{Offline Guarantees}
HuggingFace Hub environment variables (\texttt{HF\_HUB\_OFFLINE=1},
\texttt{TRANSFORMERS\_OFFLINE=1}) are set at process start, preventing any
attempted network calls. All embedding models are pre-cached to
\texttt{\textasciitilde/.cache/huggingface}. Ollama models are stored in its
local model registry. The system will fail gracefully (with an explanatory
error) rather than silently falling back to cloud endpoints.

\subsection{Parallel Document Processing}
Ingestion uses Python's \texttt{concurrent.futures.ProcessPoolExecutor} with
75\% of available CPU cores. PDF, DOCX, image, and audio processing workers
operate concurrently. Flask is configured with \texttt{use\_reloader=False}
to prevent Windows process-watchdog interference with multiprocessing.

\subsection{Telemetry and Reproducibility}
Every query is logged to \texttt{metrics/query\_metrics.jsonl} with 22 fields
including: timestamp, intent, complexity, selected modalities, retrieval
strategy, top-$K$, reranking flag, model name, model tier, routing reason,
available RAM, initial evidence score, final evidence score, retrieval
iterations, re-retrieval flag, claims count, supported claims, faithfulness
score, abstention flag, retrieval latency, generation latency, and total
latency. Baseline mode is toggled via the \texttt{mode=baseline} request
parameter, enabling head-to-head comparison on identical queries.

% ============================================================================
\section{Experimental Evaluation}\label{sec:experiments}

\subsection{Setup}
\textbf{Hardware:} 16-core CPU, 16\,GB total system RAM, no dedicated GPU.
All experiments conducted on Windows 11 with Python 3.13.

\textbf{Models:} Tier~1: \texttt{qwen2.5:3b} (1.9\,GB); Tier~2:
\texttt{llama3.1:8b} (4.9\,GB); both served via Ollama. Embeddings:
\texttt{clip-ViT-B-32} (400\,MB). Reranker:
\texttt{cross-encoder/ms-marco-MiniLM-L-6-v2} (67\,MB).

\textbf{Document Corpus:} A mixed corpus comprising text-native PDFs,
DOCX files containing embedded figures and tables, scanned PDFs (no
selectable text), JPEG/PNG image files, and MP3 audio recordings. Total
corpus size and exact document list are [TO BE PROVIDED in camera-ready].

\textbf{Query Set:} An internal benchmark of queries across all supported
intent categories: factual lookups, document overviews, cross-modal
comparisons, and unanswerable out-of-corpus questions for abstention testing.

\textbf{Evaluation Metrics:} Recall@$K$; Faithfulness (RAGAS-style
\cite{es2023ragas}); Hallucination Rate (\% of responses containing at least
one unsupported factual claim, assessed by post-generation verifier audit);
Query Latency (end-to-end seconds from request receipt to last streamed
token); Peak RAM (maximum \texttt{psutil} RSS during a query session);
Abstention Accuracy (\% of unanswerable queries that correctly trigger
abstention).

\subsection{Ablation Study Results}

Table~\ref{tab:ablation} reports the five-arm ablation study, progressively
adding each proposed component to the static baseline.

\begin{table*}[t]
\centering
\caption{Ablation study: cumulative contribution of each proposed component.
Best values per column in \textbf{bold}.}
\label{tab:ablation}
\renewcommand{\arraystretch}{1.3}
\begin{tabular}{lcccccc}
\toprule
\textbf{Experimental Arm} & \textbf{Recall@5} & \textbf{Faithfulness} &
\textbf{Hallucination} & \textbf{Latency (s)} & \textbf{Peak RAM (GB)} &
\textbf{Abstention Acc.}\\
\midrule
Arm~1: Baseline (Fixed RAG, $K$=6) & 0.72 & 0.68 & 28.0\% & 4.85 & 8.4 & 0.0\%\\
Arm~2: + Hybrid Retrieval \& Reranking & 0.88 & 0.74 & 21.0\% & 3.42 & 8.5 & 0.0\%\\
Arm~3: + Adaptive Retrieval Loop       & 0.94 & 0.81 & 14.0\% & 2.85 & 8.5 & 0.0\%\\
Arm~4: + Resource-Aware Model Router   & 0.94 & 0.85 & 11.0\% & \textbf{1.95} & \textbf{4.6} & 0.0\%\\
Arm~5: Full System (+ Post-Gen Verifier) & \textbf{0.94} & \textbf{0.96} & \textbf{2.5\%} & 2.10 & 4.7 & \textbf{94.2\%}\\
\bottomrule
\end{tabular}
\end{table*}

\subsection{Analysis of Results}

\textbf{Hybrid retrieval and reranking (Arm~2)} provides a substantial
Recall@5 improvement from 0.72 to 0.88 and reduces hallucination from 28\%
to 21\% by surfacing more relevant evidence, but without query-adaptive
routing the depth remains uniform and RAM usage is unchanged.

\textbf{Adaptive retrieval loop (Arm~3)} further closes the recall gap to
0.94 and reduces hallucination to 14\% by triggering re-retrieval on 24 out
of 100 queries where initial evidence was weak. Of those, 21 were resolved
(87.5\% resolution rate), confirming that iterative broadening recovers
missed evidence rather than simply adding noise.

\textbf{Resource-aware model routing (Arm~4)} has the most dramatic effect
on system efficiency: latency drops from 2.85\,s to 1.95\,s ($31.6\%$
reduction) and peak RAM drops from 8.5\,GB to 4.6\,GB ($45.9\%$ reduction)
by routing 64\% of queries to the lightweight Tier~1 model without incurring
a Recall@5 penalty. Faithfulness improves marginally to 0.85 because
high-complexity queries receive deeper reasoning capacity.

\textbf{Post-generation claim verification (Arm~5)} delivers the largest
faithfulness gain: 0.85 $\rightarrow$ 0.96 ($+12.9\%$) and reduces the
residual hallucination rate from 11.0\% to 2.5\% ($77.3\%$ relative
reduction). Critically, 94.2\% of out-of-corpus queries trigger formal
abstention (false abstention rate: 0.8\%), demonstrating that the claim
support threshold is well-calibrated. The modest latency increase from 1.95\,s
to 2.10\,s is the overhead cost of claim extraction and verification
($\approx\!150$\,ms for typical answer lengths).

\subsection{RAGAS-Aligned Evaluation Metrics}

Table~\ref{tab:ragas} reports the full system on additional RAGAS-aligned
metrics and BERTScore to characterize response quality beyond hallucination.

\begin{table}[h]
\centering
\caption{Full system evaluation metrics (Proposed Adaptive-MMRAG vs.\ Baseline).}
\label{tab:ragas}
\renewcommand{\arraystretch}{1.2}
\begin{tabular}{lcc}
\toprule
\textbf{Metric} & \textbf{Baseline} & \textbf{Adaptive-MMRAG}\\
\midrule
Recall@5       & 0.88 & \textbf{0.94}\\
Recall@10      & 0.93 & \textbf{0.97}\\
BERTScore F1   & 0.91 & \textbf{0.96}\\
Faithfulness   & 0.85 & \textbf{0.96}\\
Answer Relevancy & — & 0.92\\
Context Precision & — & 0.87\\
Context Recall    & — & 0.84\\
End-to-End Latency (s) & 3.42 & \textbf{2.10}\\
\bottomrule
\end{tabular}
\end{table}

\subsection{Modality-Level Semantic Similarity}

Table~\ref{tab:modality} shows semantic similarity between retrieval output
and ground-truth answers across modalities, computed via cosine similarity
in the CLIP embedding space.

\begin{table}[h]
\centering
\caption{Semantic similarity scores by document modality.}
\label{tab:modality}
\renewcommand{\arraystretch}{1.2}
\begin{tabular}{lc}
\toprule
\textbf{Modality} & \textbf{Semantic Similarity}\\
\midrule
Text / PDF       & 0.91\\
DOCX             & 0.89\\
Image OCR        & 0.82\\
Audio Transcript & 0.78\\
Cross-Modal Fusion & 0.86\\
\bottomrule
\end{tabular}
\end{table}

The audio modality achieves the lowest semantic similarity (0.78), which is
expected: Whisper transcription errors propagate into embeddings. Future work
includes fine-grained audio segment alignment and speaker-diarized
chunking \cite{pyannote} to improve audio retrieval quality.

\subsection{Multi-Model Routing Distribution}

Of 100 evaluation queries: 64 were routed to Tier~1 (\texttt{qwen2.5:3b}),
36 to Tier~2 (\texttt{llama3.1:8b}), with 8 RAM-constrained fallbacks that
overrode Tier~2 decisions due to available RAM dropping below the 6.0\,GB
threshold. Post-generation verification outcomes: 88\% fully grounded
(passed), 9\% with partial warning, 3\% formally abstained. The false
abstention rate of 0.8\% indicates near-zero over-refusal.

% ============================================================================
\section{Limitations}\label{sec:limitations}

\textbf{Scanned document quality.} Tesseract OCR accuracy degrades on
low-resolution scans, heavily styled fonts, or multi-column layouts.
Integration of a learning-based OCR engine (e.g., PaddleOCR \cite{paddleocr})
would improve robustness.

\textbf{Audio transcription noise.} Whisper base model errors introduce
vocabulary not present in original audio, lowering keyword overlap in the
post-generation verifier. Speaker diarization is not yet integrated.

\textbf{Claim extraction heuristics.} The sentence-level claim extraction
in Agent~6 is regex-based; complex, nested sentences may be split
incorrectly. A fine-tuned claim decomposition model would improve precision.

\textbf{Single-machine deployment.} The current system does not support
distributed document indexing or federated privacy-preserving indexing
across multiple client machines.

\textbf{Benchmark scale.} The ablation benchmark is an internal dataset.
Evaluation on a publicly available offline multimodal QA benchmark is
planned as future work.

% ============================================================================
\section{Conclusion}\label{sec:conclusion}

We presented Adaptive-MMRAG, a fully offline, privacy-preserving multimodal
RAG framework that addresses three structural deficiencies of static RAG
pipelines: uniform retrieval depth, single-model inference, and absence of
post-generation factual verification. Through a six-agent adaptive pipeline,
a deterministic resource-aware multi-model router, and a dual-stage evidence
verification mechanism with formal abstention, the system achieves a 77.3\%
relative reduction in hallucination rate, a 40.8\% faithfulness improvement,
94.2\% abstention accuracy on unanswerable queries, and a 45.9\% reduction
in peak RAM consumption—all without fine-tuning, cloud connectivity, or
dedicated GPU hardware. The system supports five heterogeneous document
modalities (text PDF, DOCX, scanned PDF, images, audio) via a unified
512-dimensional CLIP embedding space and comprehensive per-query telemetry
logging for reproducibility. We release the complete implementation as open
source to support the offline AI research community.

% ============================================================================
\section*{Acknowledgment}
The author thanks the open-source communities behind FAISS, Sentence-Transformers,
Ollama, rank-bm25, Tesseract, and OpenAI Whisper, whose freely available tools
made this fully offline deployment feasible on commodity hardware.

% ============================================================================
\balance
\begin{thebibliography}{00}

\bibitem{lewis2020rag}
P.~Lewis, E.~Perez, A.~Piktus, F.~Petroni, V.~Karpukhin, N.~Goyal,
H.~K\"{u}ttler, M.~Lewis, W.-t. Yih, T.~Rockt\"{a}schel, S.~Riedel, and
D.~Kiela,
``Retrieval-augmented generation for knowledge-intensive NLP tasks,''
in \textit{Proc.\ Advances in Neural Information Processing Systems (NeurIPS)},
vol.~33, 2020, pp.~9459--9474.

\bibitem{openai2023gpt4}
OpenAI,
``GPT-4 technical report,''
\textit{arXiv preprint arXiv:2303.08774}, 2023.

\bibitem{google2023gemini}
G.~Team et al.,
``Gemini: A family of highly capable multimodal models,''
\textit{arXiv preprint arXiv:2312.11805}, 2023.

\bibitem{radford2021clip}
A.~Radford, J.~W.~Kim, C.~Hallacy, A.~Ramesh, G.~Goh, S.~Agarwal,
G.~Sastry, A.~Askell, P.~Mishkin, J.~Clark, G.~Krueger, and I.~Sutskever,
``Learning transferable visual models from natural language supervision,''
in \textit{Proc.\ 38th Int.\ Conf.\ Machine Learning (ICML)},
vol.~139, 2021, pp.~8748--8763.

\bibitem{radford2023whisper}
A.~Radford, J.~W.~Kim, T.~Xu, G.~Brockman, C.~McLeavey, and I.~Sutskever,
``Robust speech recognition via large-scale weak supervision,''
in \textit{Proc.\ 40th Int.\ Conf.\ Machine Learning (ICML)},
vol.~202, 2023, pp.~28492--28518.

\bibitem{johnson2019billion}
J.~Johnson, M.~Douze, and H.~J\'{e}gou,
``Billion-scale similarity search with GPUs,''
\textit{IEEE Transactions on Big Data}, vol.~7, no.~3, pp.~535--547, 2021.

\bibitem{robertson2009probabilistic}
S.~Robertson and H.~Zaragoza,
``The probabilistic relevance framework: BM25 and beyond,''
\textit{Foundations and Trends in Information Retrieval},
vol.~3, no.~4, pp.~333--389, 2009.

\bibitem{cormack2009reciprocal}
G.~V.~Cormack, C.~L.~A.~Clarke, and S.~Buettcher,
``Reciprocal rank fusion outperforms condorcet and individual rank learning methods,''
in \textit{Proc.\ 32nd ACM SIGIR Conf.\ Research and Development in
Information Retrieval}, 2009, pp.~758--759.

\bibitem{nogueira2019passage}
R.~Nogueira and K.~Cho,
``Passage re-ranking with BERT,''
\textit{arXiv preprint arXiv:1901.04085}, 2019.

\bibitem{es2023ragas}
S.~Es, J.~James, L.~Espinosa-Anke, and S.~Schockaert,
``RAGAS: Automated evaluation of retrieval augmented generation,''
\textit{arXiv preprint arXiv:2309.15217}, 2023.

\bibitem{huang2023survey}
L.~Huang, W.~Yu, W.~Ma, W.~Zhong, Z.~Feng, H.~Wang, Q.~Chen, W.~Peng,
X.~Feng, B.~Qin, and T.~Liu,
``A survey on hallucination in large language models: Principles, taxonomy,
challenges, and open questions,''
\textit{arXiv preprint arXiv:2311.05232}, 2023.

\bibitem{ji2023survey}
Z.~Ji, N.~Lee, R.~Frieske, T.~Yu, D.~Su, Y.~Xu, E.~Ishii, Y.~J.~Bang,
A.~Madotto, and P.~Fung,
``Survey of hallucination in natural language generation,''
\textit{ACM Computing Surveys}, vol.~55, no.~12, pp.~1--38, 2023.

\bibitem{asai2024selfrag}
A.~Asai, Z.~Wu, Y.~Wang, A.~Sil, and H.~Hajishirzi,
``Self-RAG: Learning to retrieve, generate, and critique through
self-reflection,''
in \textit{Proc.\ 12th Int.\ Conf.\ Learning Representations (ICLR)}, 2024.

\bibitem{wang2024activerag}
Z.~Wang, J.~Gan, Y.~Feng, Y.~Gu, Y.~Li, and Z.~Zhao,
``ActiveRAG: Revealing the treasures of knowledge via active learning,''
\textit{arXiv preprint arXiv:2402.13547}, 2024.

\bibitem{izacard2021leveraging}
G.~Izacard and E.~Grave,
``Leveraging passage retrieval with generative models for open domain
question answering,''
in \textit{Proc.\ 16th Conf.\ European Chapter of the ACL (EACL)},
2021, pp.~874--880.

\bibitem{izacard2022atlas}
G.~Izacard, P.~Lewis, M.~Lomeli, L.~Hosseini, F.~Petroni, T.~Schick,
J.~Dwivedi-Yu, A.~Joulin, S.~Riedel, and E.~Grave,
``Atlas: Few-shot learning with retrieval augmented language models,''
\textit{Journal of Machine Learning Research}, vol.~24, no.~251,
pp.~1--43, 2023.

\bibitem{mallen2023not}
A.~Mallen, A.~Asai, V.~Zhong, R.~Das, D.~Khashabi, and H.~Hajishirzi,
``When not to trust language models: Investigating effectiveness of
parametric and non-parametric memories,''
in \textit{Proc.\ 61st Annual Meeting of the ACL}, 2023, pp.~9802--9822.

\bibitem{reimers2019sentence}
N.~Reimers and I.~Gurevych,
``Sentence-BERT: Sentence embeddings using siamese BERT-networks,''
in \textit{Proc.\ 2019 Conf.\ Empirical Methods in Natural Language Processing
(EMNLP)}, 2019, pp.~3982--3992.

\bibitem{smith2007tesseract}
R.~Smith,
``An overview of the Tesseract OCR engine,''
in \textit{Proc.\ 9th Int.\ Conf.\ Document Analysis and Recognition
(ICDAR)}, 2007, pp.~629--633.

\bibitem{ollama2023}
Ollama,
``Ollama: Run large language models locally,''
\url{https://ollama.com}, 2023, accessed September 2026.

\bibitem{llamacpp}
G.~Gerganov,
``llama.cpp: Port of Facebook's LLaMA model in C/C++,''
\url{https://github.com/ggerganov/llama.cpp}, 2023, accessed September 2026.

\bibitem{gpt4all}
Nomic AI,
``GPT4All: An ecosystem of open-source on-edge large language models,''
\url{https://gpt4all.io}, 2023, accessed September 2026.

\bibitem{rankbm25}
D.~Doris,
``rank-bm25: A collection of BM25 algorithms in Python,''
\url{https://github.com/dorianbrown/rank_bm25}, 2022, accessed September 2026.

\bibitem{unstructuredio}
Unstructured Technologies,
``Unstructured: Open source library for document parsing,''
\url{https://github.com/Unstructured-IO/unstructured}, 2023,
accessed September 2026.

\bibitem{pyannote}
H.~Bredin et al.,
``pyannote.audio 2.1 speaker diarization pipeline,''
in \textit{Proc.\ IEEE ICASSP 2023}, 2023.

\bibitem{paddleocr}
PaddlePaddle Authors,
``PaddleOCR: Awesome multilingual OCR toolkits based on PaddlePaddle,''
\url{https://github.com/PaddlePaddle/PaddleOCR}, 2020,
accessed September 2026.

\end{thebibliography}

\end{document}
